% 1. Preamble
\documentclass[12pt, a4paper]{article}

% Import Packages
\usepackage[utf8]{inputenc}        % Encoding
\usepackage{amsmath, amssymb}     % Advanced math formulas
\usepackage{mathtools}           % Math enhancements
\usepackage{graphicx}            % Image insertion
\usepackage{hyperref}            % Interactive hyperlinks & PDF bookmarks
\usepackage{color}               % Color text and backgrounds
\usepackage{geometry}            % Page geometry
\usepackage{comment}             % Multi-line comments
\usepackage{titling}             % Custom title formatting

\geometry{a4paper, margin=1in}

% Custom Command Definitions
\newcommand{\Sball}{\mathcal{S}}
\providecommand{\norm}[1]{\left\lVert#1\right\rVert}

% Hyperref Setup for Clean PDF Links
\hypersetup{
    colorlinks=true,
    linkcolor=blue,
    urlcolor=blue,
    pdftitle={Pattern Arithmetic: Verdicts from Bowls and Paths},
    pdfauthor={Sandeep Lakshminarayan Chiluveru}
}

% Metadata & Custom Title Banner
\pretitle{\begin{center}\Huge\fontfamily{tt}\selectfont ================================================================================\\[0.5em]}
\posttitle{\\[0.5em] ================================================================================\end{center}}

\title{Pattern Arithmetic: Verdicts from Bowls and Paths}
\author{Sandeep Lakshminarayan Chiluveru}
\date{\today}

\begin{document}

\maketitle

% ABSTRACT
\begin{abstract}
A companion paper, \emph{Pattern Arithmetic: Typed Algebra of Bounded
Units on a Phase-Colored Sphere}, defines a unit of truth
$x=(r,\hat n,\chi,\tau)$ on a sphere $\mathcal{S}$ and an algebra over
such units --- gather ($\oplus$), remove ($\ominus$), mix ($\boxplus$),
walk ($T$), and offset --- together with an admissibility test that
decides which sorts $\tau$ may use which operators truthfully. That
paper builds the algebra and stops at the reading: a bowl of units may
be mixed or walked, and the result is a state. It does not say how a
state, or a bowl, becomes a \emph{decision}.

This paper adds the missing step. A shaman, a tribal meteorologist, or
a market analyst holds a handful of signs and must arrive at one of a
small number of conclusions --- water or no water, storm or calm, rise
or fall. We show that the algebra already contains two distinct
mechanisms for this, corresponding to the two readings the first paper
leaves undemonstrated: an order-free reading, in which evidence
accumulates as phasors and a balanced case returns no answer rather
than a forced one; and an order-dependent reading, in which evidence
is walked as a sequence of actors to a landing compared against
seated names. We give a rule --- the \emph{swap test} --- for which
reading a given body of evidence calls for, and identify four
components a decision needs beyond what the algebra supplies: a
comparison rule, a selection rule for which signs enter and in what
order, a principled abstention, and a place for a lexicon of
conclusions to be learned rather than declared. This paper is a
framework: it specifies these four components precisely enough to be
filled in, and works one example --- the diviner's shells --- by hand.
It does not report an evaluation against data, and says exactly where
that evaluation would have to be made.
\end{abstract}

\section{Introduction}
\label{sec:intro}

A shaman throws ten shells and reads a future from them. He may throw
twenty and read it better, or forty better still --- but unless he
says so, no one else knows what he saw. The shells are evidence; the
saying is the verdict. Pattern Arithmetic gives a precise account of
the shells. It has none, so far, of the saying.

This is deliberate on the first paper's part, not an oversight. That
paper's discipline was to make the algebra trustworthy before asking
it to decide anything: an admissibility test that refuses sorts with
no honest map onto the operators, a type table that says which
combinations are legal, and a running insistence that a walk's landing
is compared to a name, never confused with one. Asking that paper to
also commit to a decision rule would have asked it to defend a claim
it was not ready to make.

It is ready now. This paper takes the atom, the bowl, and the path as
given --- cited, not re-derived --- and asks the next question: given
a bowl of evidence, what is the answer?

\subsection{A human can hold a handful; a machine need not stop
there}
A diviner's working memory bounds how many shells, how many birds, how
many signs can be weighed against one question at once. Nothing in
the algebra of the first paper imposes that bound. The claim of this
paper is narrower than it might sound: not that more evidence is
always better, but that \emph{the mechanism for combining evidence
does not change shape as the amount of evidence grows} --- the same
gather, the same two readings, apply to ten shells or ten thousand
sensor readings. Where that mechanism strains under scale (\S 4.3) is
stated plainly rather than smoothed over.

\section{What Is Inherited}
\label{sec:inherited}

We use, without re-deriving, the following from the companion paper.
Equation and axiom numbers refer to that paper.

\begin{itemize}
\item The atom $x=(r,\hat n,\chi,\tau)\in\mathcal{S}$, with $r\le 1$ a
      normalized strength, $\hat n\in S^2$ a direction, $\chi\in[0,2\pi)$
      a phase, and $\tau$ a sort tag that gates which operators may be
      truthfully applied (Axiom 1; Axiom 0).
\item Gather, $\oplus$: union on finite subsets of $\mathcal{S}$,
      sort-blind, idempotent, lossless, undone by $\ominus$ (Axiom 4).
\item Mix, $\boxplus$: a phasor reading of a bowl whose members share a
      direction and a sort,
      $\boxplus(B)=\bigl(\min(1,|z_B|),\hat n,\arg z_B,\tau\bigr)$
      with $z_B=\sum_{x\in B} r_x e^{i\chi_x}$; commutative,
      non-associative, lossy, uninvertible (Axiom 5).
\item The walk, $T$: an ordered composition
      $T_{w_n}\circ\cdots\circ T_{w_1}(x)$ of actor-parameterized
      rotations and phase shifts; non-commutative in general, and
      many-to-one, so a landing alone does not determine the walk that
      reached it (Axiom 6; \S3).
\item \emph{Levels are closed}: a unit carries no representation of
      what it constitutes, and a bowl carries no representation of the
      reading that will be applied to it. Whatever a lower unit would
      need to know about a higher one --- which units were taken as
      co-present, over what interval, under what question --- is a
      constituting fact of the higher level, not a coordinate of the
      atom (companion paper, \S2.2).
\end{itemize}

Two consequences of level closure matter enough here to restate. First,
a bowl that saturates under Mix is often a level that has been
flattened --- repeated samples of one unit, mistaken for many units at
one level (companion paper, \S2.7). Second, the atom stays a
four-tuple; grouping criteria such as a time window or a spatial
radius live outside it, deciding bowl membership before the algebra
runs rather than riding inside any single unit.

\section{Two Readings of One Bowl}
\label{sec:two_readings}

Given a bowl $B=\{x_1,\ldots,x_n\}$ of evidence bearing on one
question, the companion paper currently defines two readings that turn
$B$ into a single state: phasor Mix and ordered transformation. This is
not a claim that no other reading could be defined; it is the pair
that paper supplies today, and the pair this paper uses. We name them
here because the first paper defines both and demonstrates only one.

\subsection{Reading B: the walk}
Seat the question as a state $q$. Apply the members of $B$ as actors,
in some order $\pi$:
\[
    \mathrm{land}_\pi(q) \;=\; T_{x_{\pi(n)}}\circ\cdots\circ
    T_{x_{\pi(1)}}(q).
\]
The verdict is the landing, compared against seated names (the storm's
$\mathrm{calm}$ and $\mathrm{storm}$ pins are an instance of this).
Order matters by construction: $T_a T_b\neq T_b T_a$ in general, so
$\mathrm{land}_\pi \neq \mathrm{land}_{\pi'}$ for $\pi\neq\pi'$ is the
expected case, not a failure. This reading is fully worked in the
companion paper's storm example, which we do not repeat.

\subsection{Reading A: the phasor sum}
\label{subsec:reading_a}
Rotate every member of $B$ onto the question's axis $\hat n_q$ (a
rotation, which preserves $r$ and $\chi$), then read the whole bowl at
once:
\[
    \boxplus(B) \;=\; \Bigl(\min(1,|z_B|),\ \hat n_q,\ \arg z_B,\ \tau\Bigr),
    \qquad
    z_B=\sum_{x\in B} r_x e^{i\chi_x}.
\]
Order plays no role: $z_B$ is a sum, and sums do not care about the
order of their terms. The verdict's \emph{confidence} is $|z_B|$'s
image under the cap, and its \emph{direction} is $\arg z_B$, read
against a lexicon of conclusions seated at particular phases (e.g.\
$\chi=0$ for ``water,'' $\chi=\pi$ for ``no water''). Agreeing evidence
lengthens the sum; opposing evidence shortens it; evenly opposed
evidence returns $z_B=0$ --- black, and a principled abstention rather
than a forced choice (companion paper, Axiom 2).

\subsection{The swap test}
\label{subsec:swap_test}
Which reading applies to a given $B$ is decided by one question: does
exchanging the order in which two members of $B$ arrived change the
answer?
\begin{itemize}
\item If \textbf{no} --- the shells thrown together, ten readings
      taken at one instant --- $B$ has no order to respect, and
      Reading A is the correct mechanism. Applying Reading B here
      manufactures an order-dependence that is not in the evidence.
\item If \textbf{yes} --- birds fleeing, then a wind shift, then a
      temperature drop, where the same three signs in a different
      sequence mean something else --- $B$'s order is itself
      information, and Reading A would discard it. Reading B is
      required.
\end{itemize}
The two readings are not competing hypotheses about one mechanism;
they are the correct tool for two different shapes of evidence, and
the swap test is how a system --- human or otherwise --- tells which
shape it is holding.

\subsection{The readings compose}
\label{subsec:readings_compose}
Nothing prevents using both on one episode. A morning's observations
may arrive in a meaningful sequence (Reading B over the morning) and
also include several simultaneous signs at one moment within it
(Reading A over that moment's bowl, itself one actor in the larger
walk). This is level closure in practice (\S\ref{sec:inherited}): a
Reading A verdict is a state, and a state may act as $w$ in a $T$-path,
exactly as any other seated unit may.

\section{Four Components the Algebra Does Not Supply}
\label{sec:four_gaps}

Both readings above take a bowl and return a state. Neither, on its
own, returns a decision a person or a system can act on. Four
components stand between a state and a verdict, and none of them is
supplied by $\oplus$, $\boxplus$, $T$, or offset.

\subsection{A comparison rule}
\label{subsec:comparison_rule}
The companion paper's storm section says a ball falls \emph{near} the
$\mathrm{calm}$ pin or the $\mathrm{storm}$ pin, and does not define
near. A working comparison rule must state, at minimum: which
coordinates of the landing are read, what metric is used on those
coordinates, and whether $\tau$ must match exactly or may differ.

Two states are compared only if $\tau_A=\tau_B$; $\tau$ gates
comparison the same way it gates every other operator. Given that, a
distance is taken over all three geometric coordinates, each measured
the way the companion paper already measures it elsewhere rather than
by naive coordinate subtraction:
\begin{equation}
    d(A,B)^2 \;=\;
    w_r\,(r_A-r_B)^2
    \;+\;
    w_{\hat n}\,\arccos(\hat n_A\cdot\hat n_B)^2
    \;+\;
    w_\chi\,\bigl(\min(|\chi_A-\chi_B|,\,2\pi-|\chi_A-\chi_B|)\bigr)^2 ,
    \label{eq:comparison_metric}
\end{equation}
using the geodesic distance the offset operator already relies on for
$\hat n$, and the circular distance a phase requires for $\chi$; $r$
alone is a bounded line with no wraparound and needs neither
correction. Raw coordinate subtraction on $\hat n$ or $\chi$ is
\emph{not} a safe default: two directions or phases near $0$ and near
$2\pi$ are adjacent on the sphere or the circle but appear maximally
distant under naive subtraction, and reference atoms placed near a
natural zero point --- which is common, not rare --- would be
misjudged systematically rather than occasionally. \emph{Open in this
paper:} the weights $w_r,w_{\hat n},w_\chi$, meaning how much a
mismatch in strength, direction, or phase should count against an
otherwise close match. \S\ref{sec:worked_example} does not need this
choice made in general: every shell and every reference point there
already shares $\hat n_q$ by construction, so $w_{\hat n}$ never
enters, and the verdict is read from $\chi$ (which reference point is
nearest) with $r$ supplying the null rather than a weighted distance.
A question with reference atoms at different directions would need
$w_{\hat n}$ and $w_\chi$ actually chosen against each other, which
remains open.

\subsection{A selection rule}
\label{subsec:selection_rule}
For Reading B, the order $\pi$ that the evidence is applied in is
handed to the walk; it is not the walk's to choose. A real system must
decide which signs enter $B$ at all (\S\ref{subsec:comparison_rule}'s
window and radius criteria, companion paper \S2.2) and, when order
matters, in what sequence. This selection is itself where most of an
inference system's actual difficulty lives, and it is entirely outside
the algebra. \emph{Open in this paper.}

\subsection{A principled null}
\label{subsec:the_null}
Reading A supplies one for free: $z_B=0$. Reading B does not. A walk
always lands somewhere, and a comparison rule (\S\ref{subsec:comparison_rule})
always returns a nearest pin, however far away it is. A Reading-B
system therefore needs an explicit distance threshold beyond which the
verdict is \emph{abstain} rather than \emph{nearest pin regardless}.
This is a second place, distinct from \S\ref{subsec:comparison_rule},
where a number must be chosen rather than derived.

\subsection{Learned pins and learned signs}
\label{subsec:learned_pins}
The companion paper notes that a lookup $E_L$ may be learned
(\S5.6) without the units it points into being anything other than
fixed. The same discipline applies to inference: what is learned is
narrow. A sign's $(r,\chi)$ for a given question --- how strongly, and
in which direction, a particular bird's flight bears on ``storm'' ---
is an $E_L$ entry, learnable from labelled cases. The seated pins
$\mathrm{storm}$, $\mathrm{calm}$ are likewise $E_L$ entries, and may
be learned rather than declared. Neither requires learning a
representation in the sense of an embedding; both are learning one
small map, per sign, per question. \emph{This is the paper's one
concrete claim about the learning loop}: it is cheaper than
representation learning because it is narrow, not because it is
avoided.

\paragraph{Building a reference atom.}
\label{par:reference_atom}
A pin need not be declared; it can be built from witnessed cases the
same way live evidence is read. For a class $c$ (one storm, one
drought), let $G_c=\{y_1,\ldots,y_m\}$ be the Reading-A composite of
each witnessed episode, seated by $E_L$ from that episode's own bowl
and already rotated onto the question's axis $\hat n_q$
(\S\ref{subsec:reading_a}). Because every $y_i$ shares $\hat
n_q$ by construction, $\boxplus(G_c)$ is defined by the
companion paper's bowl-Mix equation ($\S2.7$) without needing its
still-open case of Mix across differing directions (its Future Work);
building a reference atom this way
never depends on a question that paper leaves unresolved. The
reference atom for $c$ is $\mathrm{ref}_c=\boxplus(G_c)$.

\paragraph{Coherence, and when to split.}
The strength $r$ of $\mathrm{ref}_c$ is not only a confidence value;
it is a coherence check on $c$ itself. Members of $G_c$ that agree add
constructively; members that genuinely disagree --- because $c$ is not
one pattern but two, wearing one label --- partially or fully cancel,
exactly as opposed evidence cancels in
\S\ref{sec:worked_example}. A weak $\mathrm{ref}_c$ is therefore not
noise to smooth over; it is the mechanism reporting that $c$ should be
split. Given a coherence threshold $\theta\in(0,1)$ and a minimum
group size $k_{\min}$ below which a split is not attempted: if
$r(\mathrm{ref}_c)\ge\theta$, keep $\mathrm{ref}_c$; otherwise, if
$|G_c|\ge k_{\min}$, partition $G_c$ into sub-groups and recurse on
each, so that a single label may resolve into several reference atoms
(a slow-onset and a flash drought, say); otherwise accept
$\mathrm{ref}_c$ as weak, or fall back to comparing incoming evidence
against $G_c$'s individual members rather than their compression. How
$G_c$ is partitioned when a split is warranted is not specified here
and is left open, as is the choice of $\theta$ itself; $\theta$ must
be bounded away from $0$, since every singleton group is trivially
coherent and a threshold with no floor never stops splitting. Group
size strictly decreases at each split, so the recursion terminates
regardless of $\theta$'s exact value.

\paragraph{This is level closure, applied to references.} A class
discovered to need splitting is a level-2 object in its own right
(companion paper, \S2.2): the split is discovered from $G_c$'s own
internal agreement, not asserted in advance, and each resulting
sub-reference is exactly as usable --- comparable, walkable, further
splittable --- as the class it came from. A domain with genuinely
few witnesses (two droughts, not nine) is not a failure of the
procedure; $k_{\min}$ says plainly that there is not yet evidence to
subdivide responsibly, and the honest output is a weaker reference or
a fall back to individual comparison, not a manufactured split.

\paragraph{Illustration.} Nine witnessed episodes at $r=0.6$ each, at
$\chi=8^\circ,9^\circ,9^\circ,10^\circ,10^\circ,10^\circ,11^\circ,
11^\circ,12^\circ$, compress to
$\mathrm{ref}=(1.00,\hat n_q,10.0^\circ)$: coherent, confidently kept.
The same nine witnesses, five at $\chi=0^\circ$ and four at
$\chi=180^\circ$ --- two real sub-types wearing one label --- compress
to $r=0.60$; a perfectly even split, five and five, compresses
to $r=0.00$, the same cancellation that gives
\S\ref{sec:worked_example} its abstain, now flagging the class itself
rather than a single reading of it.

\section{Worked Example: The Diviner's Shells}
\label{sec:worked_example}

We work Reading A in full, since the companion paper never carries a
phasor reading past two states, and its property demonstration
(non-associativity, on an anonymous triple) is not a reading of a
bowl that means something.

Seat the question \emph{will it rain} at $\hat n_q$. Each shell, once
cast, is read by a diviner (or, later, by a learned $E_L$) as a state
$(r,\chi)$ at $\hat n_q$: $\chi=0$ favors rain, $\chi=\pi$ favors no
rain, and intermediate $\chi$ is a hedge. Five casts, each shell
carrying $r=0.4$ unless noted:

\begin{center}
\begin{tabular}{lcc}
\hline
Cast & reading & interpretation \\
\hline
4 shells, all favor rain      & $r=1.00$ at $\chi=0^\circ$   & confident rain \\
3 favor, 1 against            & $r=0.80$ at $\chi=0^\circ$   & rain, less confident \\
2 favor, 2 against            & $r=0.00$                     & \textbf{abstain} \\
4 favor, 2 hedge at $90^\circ$ & $r=1.00$ at $\chi\approx 26.6^\circ$ & rain, hedged direction \\
12 shells ($r=0.15$ each), all weakly favor    & $r=1.00$ at $\chi=0^\circ$   & confident rain \\
\hline
\end{tabular}
\end{center}

The third row is the case worth dwelling on. Two shells favor rain and
two oppose it, at equal strength; the phasors cancel exactly, and the
verdict is black --- not a coin flip, not a default to the majority
class, but the algebra's own statement that this evidence does not
decide the question. A system built on a forced classifier (softmax
over two classes, say) has no way to express this outcome; it must
output a confident-looking number even when the evidence is silent.
Reading A's null is not an afterthought fitted onto the algebra; it is
Axiom 2, doing here exactly what it does for a Mix reading of colour.

The fifth row shows the saturation the companion paper's \S2.7
discusses: twelve weak, agreeing signs pin at $r=1$ exactly as four
strong ones do, and the difference between ``moderately confident''
and ``overwhelmingly confident'' is lost. This is the place where
\S\ref{subsec:selection_rule}'s selection rule matters most: whether
those twelve readings are twelve units at one level, or repeated
samples of one unit that should have been reduced (companion paper,
Axiom 4) before ever reaching the bowl, is a question this paper does
not resolve and flags as the clearest candidate for the next.

\emph{Reading B on the same shells} is not worked here. Doing so
honestly requires an order for the casts, which a physical handful of
shells thrown at once does not supply --- exactly the case the swap
test (\S\ref{subsec:swap_test}) assigns to Reading A. A worked Reading
B example belongs to evidence that genuinely arrives in sequence, and
the companion paper's storm already provides one.

\section{Level Closure in Practice: Correcting a Path}
\label{sec:level_closure_practice}

A $T$-path records the sequence of actors applied to it, faithfully,
by definition (companion paper, Axiom 6). It has no means of knowing
whether that sequence was the right one --- level closure
(\S\ref{sec:inherited}) says precisely this: a lower unit carries no
representation of what it constitutes, and neither the atom nor the
path can see the higher-level context that would tell it it has erred.

A concrete case: a sequence of observed signs is walked in the order
observed, but a later sign (or a completed higher-level unit, such as
the rest of a sentence or the rest of a song) makes that order appear
wrong --- \emph{twinkle, little, twinkle, star} rather than
\emph{twinkle, twinkle, little, star}. Correcting this is a level-2
operation: it requires comparing the recorded path against context the
path itself does not have access to, and revising the walk
probabilistically. This paper's position is that the path must
therefore be \emph{legible} to the level above --- readable as an
ordered record, not only walkable --- so that a corrector at the next
level can inspect and revise it. Nothing in the companion paper's
axioms prevents this; nothing in them provides it either. We flag it
as the clearest point of contact between this paper's verdicts and a
learning loop operating above them, and leave the corrector itself
unspecified.

\section{Scope and Open Questions}
\label{sec:scope}

This paper is a framework, in the same sense the companion paper is:
it specifies mechanism and states honestly what it has not
demonstrated at scale. \emph{[OPEN --- decide before circulating: does
this paper stop here, as a specification plus one worked toy example,
or does it commit to an evaluation against real data (e.g.\ a labelled
forecasting or forex task) with a stated baseline? The two choices
imply different remaining sections and different claims in the
abstract; see the companion paper's own title history for how much
that choice shapes what a title may honestly claim.]}

Regardless of that choice, the following are stated open rather than
resolved by this paper, and are drawn directly from
\S\ref{sec:four_gaps}:

\begin{itemize}
\item The comparison rule (\S\ref{subsec:comparison_rule}) is a
      choice, not a derivation; a different metric may be more honest
      for a given domain.
\item The selection rule (\S\ref{subsec:selection_rule}), including
      how the grouping window of companion-paper \S2.2 is set, is the
      largest source of real difficulty and is entirely unaddressed
      here.
\item The saturation--under--aggregation problem
      (\S\ref{sec:worked_example}, companion paper \S2.7 and \S2.3) is
      diagnosed, not solved: whether to route repeated samples through
      a reduction before the bowl, or revisit the companion paper's
      Axiom~1 note on an order-preserving bijective bound, is left to
      whichever setting first needs to aggregate at scale.
\item The path corrector (\S\ref{sec:level_closure_practice}) is named
      and not built.
\end{itemize}

\end{document}
